%! Functions and Function Notation — Unit 1, Lesson 1.1 — Lesson Plan
% Section order per LESSON_SHAPE.md §5. 1.1 now OPENS Unit 1 (the old 1.0 was retired), so no
% box here refers back to a previous lesson. The evaluate-vs-solve "two directions" material
% that used to be this lesson's crux has moved to 1.4 (Inverse Functions), where f(a)=b <=>
% f^{-1}(b)=a is the same idea in its natural home; the crux here is the piecewise boundary.
\documentclass[10pt]{article}
\usepackage{precalculus-article}
\usepackage{precalculus-boxes}
\usepackage{graphicx}
\graphicspath{{images/}}
\newcommand{\TallMath}[1]{$\displaystyle #1\rule[-1.4em]{0pt}{3.2em}$}
\newcommand{\CourseName}{Precalculus}
\newcommand{\MeetingLength}{60 minutes}
\newcommand{\UnitNumberName}{Unit 1: Functions and Change \quad}
\newcommand{\LessonNumberName}{Lesson 1.1: Functions and Function Notation}

\begin{document}
\begin{center}
    {\Large\bfseries \CourseName \\
    {\small\bfseries \UnitNumberName \LessonNumberName}}
\end{center}

\begin{tcolorbox}[colback=lilac, colframe=plum, boxrule=0.9pt, arc=2mm,
    left=2mm, right=2mm, top=1.5mm, bottom=1.5mm]
    \textbf{Primary Objective:} Students will read and write function notation
    --- $f(x)$ --- evaluate a function given by a formula, a table, or a graph,
    state the domain and range of a function presented as a list of points or a
    graph, and evaluate a \textit{piecewise} function by deciding which rule owns
    the input --- including when the input sits exactly on a boundary.
    \hfill\textit{\small (Foundations ramp --- no CED topic)}
\end{tcolorbox}

\boxguard[22]
\begin{skillbox}[Lesson Flow --- Gradual Release (60 minutes)]{lilac}
\small
\renewcommand{\arraystretch}{1.25}
\begin{tabularx}{\linewidth}{@{} l c X @{}}
\textbf{Phase} & \textbf{Min} & \textbf{What students are doing, and where it lives} \\
\midrule
Warm-Up                       & 5  & Silent, individual spiral review --- warm-up sheet. \\
Hook                          & 2  & Whole-class discussion --- hook box atop the guided notes. \\
\textbf{I Do} --- model       & 10 & Watch and annotate while the teacher thinks aloud --- the NOTATION and EVALUATING rows: definitions, displays, problems 1 and 3. \\
\textbf{We Do} --- guided     & 22 & Fill the notes together, every problem cold-called --- problems 2, 4--6, then the DOMAIN \& RANGE and PIECEWISE rows (7--14; the teacher works 7 and 11). \\
\textbf{You Do} --- independent & 13 & Work alone while the teacher circulates --- the You Do row, problems 15--18. \\
Debrief                       & 5  & Whole-class share-out and the headline. \\
Homework Launch               & 3  & Start homework problems 1 and 4 alone; the teacher reads 4(b) over shoulders. \\
\end{tabularx}

\smallskip
\textbf{This is the tightest period in the unit.} Four ideas --- notation,
evaluating, domain/range, and piecewise rules --- share one block, so the guided
phase runs 22 minutes. The whole release lives in the guided notes; there is no
activity sheet to raid when it overruns. If the clock slips, cut \textbf{problem 18}
(writing a tariff as a piecewise rule) first, then problem 17, and protect the
\textbf{We Do} boundary work in the PIECEWISE row: homework problems 4 and 8 cannot
be done without it. Never cut the Homework Launch to zero --- it is the only
formative read of the period now that there is no exit ticket.
\end{skillbox}

\renewcommand{\arraystretch}{1.6}
\begin{skillbox}[Priority Ideas \& Skills]{goldbox}
\begin{tabularx}{\linewidth}{@{}X X@{}}
\textbf{Priority Skills} &
\textbf{Key Understandings} \\
\midrule
Read and write function notation. &
\begin{itemize}[leftmargin=1.2em,itemsep=2pt,topsep=0pt]
  \item $f(x)$ is read ``$f$ of $x$'' --- never ``$f$ times $x$.''
  \item One statement $f(3) = 7$ carries three facts: input 3, output 7,
        and the point $(3, 7)$ on the graph.
  \item An input need not be a number --- only something the rule turns into
        exactly one output.
\end{itemize} \\
Evaluate a function from a formula, a table, or a graph. &
\begin{itemize}[leftmargin=1.2em,itemsep=2pt,topsep=0pt]
  \item The slot takes anything: $f(\square) = 5(\square) - 8$. Write the
        parentheses first, then fill them.
  \item Negative inputs are the error magnet: $(-4)^2 \ne -4^2$.
\end{itemize} \\
State the domain and range of a function. &
\begin{itemize}[leftmargin=1.2em,itemsep=2pt,topsep=0pt]
  \item Domain = every input the function uses; range = every output it
        produces (a repeat listed once).
  \item From a graph, both are read as inequalities --- $0 \le x \le 8$ --- or
        with brackets, $[\,0, 8\,]$.
\end{itemize} \\
Evaluate a piecewise function. &
\begin{itemize}[leftmargin=1.2em,itemsep=2pt,topsep=0pt]
  \item The condition beside a rule is its \textit{address}: it says which
        inputs that rule is responsible for.
  \item Decide the rule \textit{first}, then evaluate only that one. Running
        both rules and choosing is the error.
  \item \textbf{Target misconception:} at a boundary, students use the rule
        whose graph \textit{reaches} the input rather than the rule whose
        condition \textit{owns} it --- the open circle versus the filled dot.
\end{itemize} \\
\end{tabularx}
\end{skillbox}

\begin{skillbox}[{Vocabulary, Concepts, \& Theorems}]{greenbox}
\renewcommand{\arraystretch}{1.4}
\begin{tabularx}{\linewidth}{@{} >{\leavevmode\bfseries\color{plum}}p{0.24\linewidth} X @{}}
Function notation & The symbol $f(x)$: $f$ names the rule, $x$ is the input, $f(x)$ is the output at $x$. Read ``$f$ of $x$.'' \\
Evaluate          & Find the output a function produces for a given input. \\
Domain            & The set of all inputs a function uses. \\
Range             & The set of all outputs a function produces; a repeated output is listed once. \\
Piecewise function & One function built from two or more rules, each rule owning its own stretch of inputs. \\
Undefined         & What we call an expression with no value, such as a division by zero --- inputs causing it are excluded from the domain. \\
\end{tabularx}
\end{skillbox}

\begin{skillbox}[Activate Prior Knowledge \& Spiral Review (5 min)]{lilac}
    Please see attached warm-up (spiral review).
    Skills reviewed: evaluating a linear expression at a negative value,
    order of operations with a squared negative,
    solving a two-step linear equation,
    and deciding whether a number satisfies an inequality such as $x \ge 2$ ---
    the test every piecewise condition is made of.
\end{skillbox}

\begin{skillbox}[Hook (2 min)]{lilac}
    \textbf{The machine gets a name.} A vending machine you can trust: press the
    same button, get the same item, every time. Let $P(b)$ be the price in dollars
    of the item at button $b$.
    \begin{enumerate}[itemsep=2pt]
        \item On the board, write only: $P(\text{B4}) = 1.25$. Ask what the
              equation is \textit{saying}. Push for a full sentence --- ``the
              item at button B4 costs \$1.25.''
        \item Follow up: what is the input here? It is \textbf{B4} --- a button,
              not a number. Inputs do not have to be numbers; they only have to
              produce exactly one output.
    \end{enumerate}
    \textit{Discussion prompt:} one short symbol just packed in the machine, the
    input, and the output. That efficiency is the whole point of function
    notation --- and it is why every lesson from here on is written in it.
\end{skillbox}

\begin{skillbox}[Lesson --- I Do / We Do / You Do]{lilac}
\begin{multicols}{2}

\textbf{I DO --- Part 1: Notation and the Empty Slot} (10 min)

\textit{Notes rows NOTATION and EVALUATING: read each definition, mark up the
display, work problems 1 and 3. Teacher works; students watch and annotate.
Nothing is cold-called in this phase.}
\begin{itemize}
  \item Read the NOTATION sentences aloud: $f$ is the machine's name, $x$ is the
        input slot, $f(x)$ is the output. Insist on ``$f$ of $x$,'' and confront
        the misreading head-on --- those parentheses hold the input; they do not
        multiply.
  \item Mark up the machine display: $F(m) = 2.50 + 1.75m$ for a rideshare fare.
        Fill the two label boxes yourself --- $m$ is miles driven, $F(m)$ the fare
        in dollars.
  \item Work \textbf{problem 1} aloud, the three-facts reading: $f(3) = 7$ says
        input 3, output 7, point $(3,7)$.
  \item Read the EVALUATING sentences and the slot display, then work
        \textbf{problem 3} by doing the habit, not asking for it:
        $f(x) = 5x - 8$ becomes $f(\square) = 5(\square) - 8$. Write the
        parentheses first, \textit{then} fill them, narrating both moves.
\end{itemize}

\textbf{WE DO --- Part 2: Evaluating, Negatives Included} (6 min)

\textit{Problems 2 and 4--6. Class fills the blanks together; cold-call every one.}
\begin{itemize}
  \item \textbf{Problem 2} first --- back to the NOTATION grid: say what $F(4)$
        is asking for before anyone computes 9.50.
  \item \textbf{Problems 4 and 5}: $f(-2)$, then $g(x) = x^2 + 3x$ at $g(-4)$ ---
        two negatives to track in one substitution.
  \item \textbf{Problem 6} is the classic trap $-4^2$ vs.\ $(-4)^2$: make the
        error visible on the board once, on purpose, and let the class catch it.
\end{itemize}

\columnbreak

\textbf{WE DO --- Part 3: Domain and Range} (7 min)

\textit{Notes row DOMAIN \& RANGE. Class fills the blanks together; cold-call
every one. This is new vocabulary, not practice.}
\begin{itemize}
  \item Definitions first: domain = all inputs, range = all outputs. Then the
        display: label the gold bracket \textit{domain} and the red bracket
        \textit{range} on the board.
  \item Work \textbf{problem 7} yourself --- from a list of points both sets are
        just lists, with the repeated 6 written once.
  \item \textbf{Problem 8} reads the graph: $0 \le x \le 8$ and $1 \le y \le 7$,
        then $[\,0, 8\,]$ --- say plainly that both forms are acceptable all year.
  \item \textbf{Problem 9} is the formula, first look: $k(x) = 8/(x+3)$ has no
        output at $x = -3$. One example only; the general treatment belongs to
        Unit 3. \textbf{Problem 10} lands ``a repeat is listed once'' and keeps
        many-to-one alive.
\end{itemize}

\textbf{WE DO --- Part 4: Piecewise Functions} (9 min)

\textit{Notes row PIECEWISE --- the spine of the lesson. Class fills the blanks
together; cold-call every one. Teacher works problem 11; the class works 12--14.}
\begin{itemize}
  \item Read the three steps aloud and make the class say the key word:
        the condition is the rule's \textbf{address}. No input belongs to two
        rules.
  \item Mark up the graph before any computing. The open circle at $(3,6)$ is the
        whole lesson in one dot: the top rule \textit{reaches} $x = 3$ but its
        address stops at $x < 3$. The filled dot at $(3,2)$ is the output the
        function actually has. Fill the two label boxes yourself --- ``bottom''
        and ``2'' --- then give one minute for the ``in your own words'' line.
  \item Work \textbf{problem 11}, $f(1)$, naming the condition out loud before
        computing. \textbf{Problem 13 is the crux}: $f(3) = 2$, not 6. Ask for the
        \textit{reason}, not the number. \textbf{Problem 14} moves the same
        decision to a formula with no graph and a negative boundary, $x \le -1$.
\end{itemize}

\textbf{YOU DO --- Part 5: The You Do Row} (13 min)

\textit{Notes row TABLE \& FORMULA, problems 15--18 --- the whole independent
phase. Students work alone; circulate and say nothing a neighbor could say
instead.}
\begin{itemize}
  \item \textbf{Problem 15} evaluates $f(x) = 3x - 7$ at a positive and a negative
        input; \textbf{16} is the crux transferred --- $h$ has no graph, so the
        boundary $x = 2$ must be settled from the condition alone.
  \item \textbf{Problem 17} returns to a table: one lookup, then domain and range
        with the repeated 9 listed once.
  \item \textbf{Problem 18} is the finisher's question --- a garage tariff in
        words, to be \textit{written} as a piecewise rule. Constructing the
        conditions is harder than reading them, and it is what homework problem 8
        asks again.
\end{itemize}
\end{multicols}
\end{skillbox}

\begin{skillbox}[Active Monitoring]{redbox}
Circulate during the \textbf{You Do} phase --- the thirteen minutes where students
are working problems 15--18 without you --- and check. The crux is
\textbf{problem 13} ($f(3) = 2$, not 6); its You Do transfer is \textbf{problem
16} ($h(2)$, with no graph to look at).
\begin{itemize}
  \item \textbf{Common error:} reading $f(x)$ as multiplication --- watch for
        students ``solving for $f$'' or dividing both sides by $f$. Cold-call:
        ``In $F(m)$, what is the machine's name? What goes in the slot?''
  \item \textbf{Common error:} dropping parentheses on negative inputs, so
        $g(-4) = -4^2 + 3(-4)$. Ask the student to say the substitution aloud
        using the word ``quantity.''
  \item \textbf{Common error (the crux):} at $h(2)$, using the rule whose formula
        comes first, or the one whose graph reaches the boundary. Push: ``Read me
        the condition. Does 2 pass that test?'' Point at the open circle they
        annotated in Part 4.
  \item \textbf{Common error:} evaluating \emph{both} rules and picking the
        friendlier answer. Ask: ``How many outputs may one input have?''
  \item \textbf{Domain/range:} students listing a repeated output twice, or
        swapping the two sets. Cold-call: ``Domain --- inputs or outputs?''
\end{itemize}
\end{skillbox}

\begin{skillbox}[Differentiation --- During You Do (13 min)]{redbox}
There is no separate activity sheet: students work the notes' You Do row
(problems 15--18) alone while you circulate. Differentiate by \textit{where you
stand}, not by handing out different paper.

\begin{itemize}
  \item \textbf{Support.} Sit with a struggling student on \textbf{problem 15}
        and make them write $f(\square) = 3(\square) - 7$ before substituting.
        The empty-slot habit from Part 1 is the whole fix for $f(-1)$; if they
        write the substitution straight down with no parentheses, the sign error
        is already in.
  \item \textbf{On level.} \textbf{Problem 16} is the boundary check. Do not
        answer it; have them read the two conditions aloud, decide which test
        $x = 2$ passes, and then walk away.
  \item \textbf{Extend.} \textbf{Problems 17 and 18} are the finisher's
        questions. 17 rewards a student who lists the repeated output 9 only
        once. 18 asks them to \textit{write} the conditions rather than read
        them --- a student who notices the second rule must start \emph{above}
        2 hours, not at it, has understood the boundary from the other side.
        Either answer presents in the Debrief.
\end{itemize}
\end{skillbox}

\boxguard[20]
\begin{skillbox}[Debrief (5 min)]{redbox}
Whole class, teacher-led, packets still open. Three moves, in order:
\begin{enumerate}[itemsep=3pt]
  \item \textbf{Share out the You Do} (3 min) --- one answer per problem,
        not a full review. Problem 15: $f(-1)$, said aloud with the parentheses.
        Problem 16: the three values, each with the condition that chose the rule.
        Problem 17: the range with 9 listed once. Problem 18: the two conditions.
        If a student wrote 18's conditions cleanly, they present it --- writing a
        boundary is the hardest thing in the lesson, and it is the only exposure
        the rest of the class gets before homework problem 8 asks for it.
  \item \textbf{Name the headline} (1 min) --- say it, then make the class say it
        back: \textit{the condition beside a rule is its address. Decide which rule
        owns the input first; then evaluate only that one.} Point at the open
        circle on the notes graph while you say it.
  \item \textbf{Point forward} (1 min) --- 1.2 stops running one machine at a
        time and starts combining two of them into a third.
\end{enumerate}
\textbf{Do not} re-teach domain and range here, take new questions, or let the
homework start early. If the launch shows it did not land, it is the 1.2 warm-up's
job, not this minute's.
\end{skillbox}

\begin{skillbox}[Homework Launch (3 min)]{redbox}
\textbf{There is no exit ticket.} The formative read comes twice: circulating the
You Do, and again here. Write the due date on the board, have students copy it into
the cover's Homework row, then start \textbf{homework problems 1 and 4} alone and in
silence (3 minutes). Circulate straight to the students you flagged during You Do.
\begin{itemize}[itemsep=2pt]
  \item \textbf{Problem 1(a)} asks what $C(3)$ \textit{means} before anyone
        computes --- a student who writes ``$C$ times 3'' is still reading the
        parentheses as multiplication.
  \item \textbf{Problem 4(b)} is the diagnostic: $D(3)$ sits exactly on the
        boundary, and $D(3) = 4$. A student who writes 4 but cannot name the
        condition has guessed; a student who writes 13 used the rule that reaches
        rather than the rule that owns.
\end{itemize}
\textbf{Sort what you see into three piles} as you walk: \textit{boundary solid}
(nothing to do), \textit{boundary shaky} (a 1.2 warm-up item on reading a
condition), and \textit{still multiplying $f$ by $x$} (a one-minute check-in at
the door of 1.2).
\end{skillbox}

\begin{skillbox}[Reinforcement \& Extension]{goldbox}
\textbf{Homework --- printed, graded (2 pages).} Last in the packet; due on the date
the teacher sets. Every context is new. \textit{Part A} (1--3) evaluates from a
formula --- a phone plan in context, a squared negative, and $-t^2$ at a negative
input, the $-4^2$ trap in a new costume. \textit{Part B} (4--5) is piecewise: a
courier tariff evaluated on both sides of its boundary and \emph{on} it, then a
formula-only piecewise with a negative boundary. \textit{Part C} (6--7) reads a new
graph and states domain and range in both notations. \textit{Part D} (8--10) builds
a piecewise rule from a pool tariff, finds a forbidden input, and problem 10 asks
for the headline in writing.

\textbf{AP Practice --- extra credit (2 pages).} Optional, before the homework in the
packet. Section I is five AP-style multiple-choice items (four options) whose
distractors are the lesson's errors: $f(x)$ as multiplication (1), dropped
parentheses on a negative input (2), the rule that reaches the boundary instead of
the rule that owns it (3), domain and range swapped (4), and the coefficient missed
when clearing a denominator (5). Section II is one free-response question on a
shipping tariff: evaluate and interpret in context (a), why a boundary value is not
ambiguous (b), what breaks if no rule owns the boundary (c), a context domain (d),
and writing a third rule with its condition (e) --- the finisher's question.

\textbf{Preview:} Next lesson (1.2, Function Operations) stops running one machine
at a time: two functions are added, subtracted, multiplied, and divided to build a
third, and the domain of the result is decided by both parents at once. The
homework's closing box points back at the phone plan in problem 1 to set it up.
\end{skillbox}

\begin{teachernote}[Warm-Up]
Five minutes, silent start, no calculators. Items 1--3 are the exact algebra
today's evaluating runs on --- a student who misses item 1 or 2 will make the
same sign error inside $f(-2)$ within the hour, so note names. Item 4 looks
trivial and is not: it is the only rehearsal students get for reading a piecewise
condition as a pass/fail test, and the boundary value 2 appears in both parts on
purpose. A student who circles 2 in \emph{both} (a) and (b) has the exact
misconception Part 4 is built to break --- flag them before the notes start.
\end{teachernote}

\begin{teachernote}[Guided Notes]
The notes are one Main Ideas / Notes table, four instruction rows and a You Do
row; in each instruction row you read the printed definition, mark up the
display, and work the first problem of the grid (1, 3, 7, 11), and the class
works the rest with the pen in their hand, every problem cold-called. The
NOTATION row lives or dies on the language: make students \textit{say} ``$f$ of
$x$'' aloud several times, and translate $f(3) = 7$ (problem 1) into a sentence
before anyone computes. In the EVALUATING row the $\square$ display is the whole
point --- write the empty parentheses first, every time, and let problem 6 catch
the $-4^2$ trap in the open. The DOMAIN \& RANGE row comes \textit{before}
piecewise on purpose: students need ``which inputs does this function use'' as a
habit before a rule starts owning only some of them.

PIECEWISE is the spine. Do the graph before the algebra --- the open circle at
$(3,6)$ and the filled dot at $(3,2)$ are the misconception drawn, and a student
who annotates those two dots has the lesson. Then problem 13 asks for the
\emph{reason}, not the number: the top rule reaches $x = 3$ but its address stops
at $x < 3$. Do not accept ``because it says so''; make them read the condition
aloud. Problem 14 is the same decision with no picture and a negative boundary,
which is where it gets hard.

The You Do row (15--18) is the entire independent phase and carries thirteen
minutes, so treat it as the lesson's third act. Problem 16 is the crux with the
graph taken away; problem 18 is the crux inverted --- writing conditions rather
than reading them. If the clock slips, cut problem 18 and let the Debrief carry
it; it returns as homework problem 8 and AP Practice item 6(e) either way.
\end{teachernote}

\begin{teachernote}[Debrief]
Five minutes, and on this lesson they are genuinely at risk --- the guided phase
runs 22 minutes and will try to borrow them. Borrow from problem 18 instead;
it is the one piece of the period that returns in the homework (problem 8). Do
not borrow from the three-minute Homework Launch either: with no exit ticket, it
is the last look you get before the work leaves the room. The share-out is one
answer per You Do problem so that a student who stalled on 15 still hears the
boundary result in 16 --- the exact shape of homework problem 4, which they will
otherwise meet alone at home.
\end{teachernote}

\begin{teachernote}[AP Practice]
Extra credit, optional, and not a substitute for the homework --- say so when you
hand out the packet. Score it however extra credit works in your gradebook; the
cover gives it a $+$ cell outside the total. Multiple-choice answers are B, D, A, A,
C. Item 3 is the most useful one to glance at: a student who circles (B) 13 used
$3x + 1$ because it is written first, and a student who circles (D) thinks both
rules apply --- those are the two halves of the lesson's misconception, and either
one means problem 13 of the notes did not land. In the free response, accept any
correct inequality or bracket form in (d); part (c) is the deepest question on the
page --- credit any answer saying that no rule would own $w = 2$, so $C(2)$ would be
undefined and 2 would fall out of the domain.
\end{teachernote}

\begin{teachernote}[Homework]
Graded, two pages, due on the date you set; students start problems 1 and 4 in the
Homework Launch. Grade problems 4 and 10 for accuracy and the rest for completion
and shown work if time is short: 4(d) (naming the condition that owns $m = 3$) and
10 (why exactly one rule applies) are the two items that say whether the boundary
landed, and together they predict 1.2 better than the evaluating items do. On
problem 3, $q(-2) = 10$ is the $-(-2)^2 = +4$ slip --- mark it, and if it is common,
reuse problem 6 of the notes in the 1.2 warm-up. Problem 8 is the one worth reading
closely: a student who writes the second condition as $12 \le a \le 80$ has given
age 12 two prices, which is the boundary error committed from the writing side.
\end{teachernote}
\end{document}
